\chapter{Fundamentos de RAG}
\label{chap:rag_fundamentos}

\chapterepigraph{``Knowledge is indivisible.''}{Isaac Asimov, \textit{The Roving Mind} (1983).}
{\small\itshape\color{AccentGray}
A afirma\c{c}\~{a}o de Asimov sobre a indivisibilidade do conhecimento~\cite{asimov1983rovingmind} encontra eco preciso na pedagogia de Paulo Freire: nenhuma palavra faz sentido desconectada do mundo que lhe d\'{a} origem~\cite{freire1968pedagogia}. Freire formulou esse princ\'{i}pio como imperativo metodol\'{o}gico --- ler o mundo antes de ler a palavra, pois a compreens\~{a}o do contexto precede e ancora qualquer produ\c{c}\~{a}o textual. O pipeline RAG operacionaliza esse princ\'{i}pio de forma literal: antes de gerar qualquer resposta, o sistema recupera fragmentos do mundo real --- documentos, registros, hist\'{o}ricos --- e os oferece ao modelo como contexto irredut\'{i}vel. Sem recupera\c{c}\~{a}o, o modelo l\^{e} apenas palavras sem ancoragem na realidade concreta do dom\'{i}nio.\par}
\vspace{0.6em}
% ---------------------------------------------------------------
\section{Pr\'e-requisitos}
% ---------------------------------------------------------------

Para aproveitar este cap\'itulo, o leitor deve ter conclu\'ido:

\begin{itemize}[leftmargin=*]
  \item Cap\'itulo~\ref{chap:fundamentos_llm_ia} --- tokeniza\c{c}\~ao,
        janela de contexto e \textit{prompting};
  \item Cap\'itulo~\ref{chap:agentes_langgraph_memoria} --- agentes
        ReAct com LangGraph;
  \item Conceito de vetores e produto interno (\'algebra linear b\'asica).
\end{itemize}

% ---------------------------------------------------------------
\section{Escopo do cap\'itulo}
% ---------------------------------------------------------------

\begin{quote}\itshape
\textbf{Onde estamos na hist\'oria:} loop interno, mem\'oria do agente.
O agente j\'a age e usa ferramentas (caps.~3--4). Agora ele precisa
de mem\'oria de longo prazo: documentos corporativos que n\~ao cabem
no contexto do LLM. RAG \'e a ponte entre o conhecimento externo e
a janela de contexto. Os cap\'itulos~6 e~9 aprofundar\~ao esse
mecanismo com fontes heterog\^eneas e grafo de conhecimento.
\end{quote}

Este cap\'itulo cobre o pipeline b\'asico de
\textit{Retrieval-Augmented Generation}\index{RAG|textbf}
(RAG): ingest\~ao, \textit{chunking}, gera\c{c}\~ao de
\textit{embeddings}, indexa\c{c}\~ao em FAISS e recupera\c{c}\~ao
sem\^antica. O objetivo \'e construir, com c\'odigo execut\'avel
localmente e sem depend\^encia de AWS, um pipeline RAG minimal
que responde a perguntas usando documentos indexados.

% ---------------------------------------------------------------
\section{Objetivos de aprendizagem}
% ---------------------------------------------------------------

Ao concluir este cap\'itulo, o leitor ser\'a capaz de:

\begin{itemize}[leftmargin=*]
  \item Explicar por que RAG reduz alucina\c{c}\~ao em sistemas
        de IA para dom\'inios espec\'ificos;
  \item Implementar \textit{chunking} por par\'agrafos com controle
        de tamanho;
  \item Gerar e persistir \textit{embeddings} com FAISS;
  \item Executar busca sem\^antica por similaridade de cosseno;
  \item Construir um pipeline RAG completo: documento $\to$ chunks
        $\to$ vetores $\to$ busca $\to$ resposta com cita\c{c}\~ao.
\end{itemize}

% ---------------------------------------------------------------
\section{Conceitos te\'oricos}
% ---------------------------------------------------------------

\subsection{Por que RAG?}

LLMs armazenam conhecimento nos par\^ametros durante o treinamento.
Esse conhecimento \'e est\'atico: n\~ao atualiza ap\'os o \textit{cutoff}
de treinamento e n\~ao cont\'em informa\c{c}\~ao propriet\'aria da
organiza\c{c}\~ao. Em dom\'inios corporativos, a resposta correta
depende de pol\'iticas atuais, contratos e dados que mudam com
frequ\^encia --- incompat\'ivel com re-treinamento constante.

RAG resolve esse problema em dois passos: (1) no momento da
pergunta, recupera os fragmentos mais relevantes do reposit\'orio de
documentos; (2) injeta esses fragmentos no \textit{prompt} como
contexto, permitindo que o LLM gere uma resposta fundamentada.

\begin{figure}[htbp]
\centering
\resizebox{\textwidth}{!}{%
\begin{tikzpicture}[
  node distance=1.2cm and 1.8cm,
  box/.style={draw, rounded corners, text width=2.8cm, align=center,
              minimum height=0.8cm, font=\small},
  arr/.style={-Stealth, thick}
]
  \node[box, fill=blue!10]  (doc)   {Documentos\\corporativos};
  \node[box, fill=blue!10, right=of doc]  (chunk) {Chunking};
  \node[box, fill=blue!10, right=of chunk](embed) {Embeddings};
  \node[box, fill=blue!10, right=of embed](index) {\'Indice FAISS};

  \node[box, fill=green!10, below=1.5cm of doc]  (query) {Pergunta\\do usu\'ario};
  \node[box, fill=green!10, right=of query]       (qemb)  {Embedding\\da query};
  \node[box, fill=green!10, right=of qemb]        (busca) {Busca por\\similaridade};
  \node[box, fill=orange!15,right=of busca]       (ctx)   {Contexto\\recuperado};
  \node[box, fill=orange!15,below=0.8cm of ctx]   (llm)   {LLM + Prompt\\com contexto};
  \node[box, fill=orange!15,left=of llm]          (resp)  {Resposta\\com cita\c{c}\~ao};

  \draw[arr] (doc)   -- (chunk);
  \draw[arr] (chunk) -- (embed);
  \draw[arr] (embed) -- (index);
  \draw[arr] (query) -- (qemb);
  \draw[arr] (qemb)  -- (busca);
  \draw[arr] (index) -- (busca);
  \draw[arr] (busca) -- (ctx);
  \draw[arr] (ctx)   -- (llm);
  \draw[arr] (llm)   -- (resp);
\end{tikzpicture}
}
\caption{Pipeline RAG: fase de ingest\~ao (azul) e fase de recupera\c{c}\~ao (verde/laranja).}
\label{fig:pipeline_rag}
\end{figure}

\subsection{Chunking: dividir para indexar}

Um documento t\'ipico tem centenas de par\'agrafos. Indexar o
documento inteiro geraria um vetor \'unico que n\~ao captura
nuances locais. O \textit{chunking}\index{chunking|textbf} divide
o documento em fragmentos menores, cada um com sua pr\'opria
representa\c{c}\~ao vetorial.

A estrat\'egia adotada no projeto \'e \textbf{chunking por
par\'agrafos com janela deslizante de tamanho fixo}:

\begin{itemize}[leftmargin=*]
  \item Par\'agrafos s\~ao delimitados por linhas em branco;
  \item Chunks acumulam par\'agrafos at\'e atingir 400 palavras;
  \item Par\'agrafos individuais maiores que 400 palavras n\~ao
        s\~ao truncados (ficam em chunk pr\'oprio).
\end{itemize}

\subsection{Embeddings: representar significado em vetores}

Um \textit{embedding}\index{embeddings|textbf} \'e um vetor
denso de alta dimens\~ao ($d = 1536$ para os modelos da Bedrock)
que representa o significado sem\^antico de um texto. Textos
semanticamente similares geram vetores pr\'oximos no espa\c{c}o
euclidiano, o que permite busca por similaridade eficiente.

A similaridade entre dois vetores \'e medida pelo cosseno:

$$\text{cos}(\mathbf{u}, \mathbf{v}) = \frac{\mathbf{u} \cdot \mathbf{v}}{\|\mathbf{u}\| \cdot \|\mathbf{v}\|}$$

Vetores normalizados reduzem a busca de cosseno a um produto
interno, que o FAISS\index{FAISS|textbf} executa com \'otimos
de hardware (SIMD/BLAS).

\subsection{FAISS: \'indice de vetores local}

O \textit{Facebook AI Similarity Search}
(FAISS)~\cite{faiss2017} \'e uma biblioteca de indexa\c{c}\~ao
de vetores que funciona inteiramente em mem\'oria, sem
necessidade de servidor externo. Para o loop interno
(desenvolvimento offline), o FAISS \'e a op\c{c}\~ao natural:
r\'apido, reproduz\'ivel e portável. Para o loop externo (produ\c{c}\~ao
com concorr\^encia), o projeto migra para OpenSearch\index{OpenSearch}
(Cap.~\ref{chap:rag_corporativo}).

O \'indice \texttt{IndexFlatL2} armazena vetores sem compress\~ao e
executa busca exata por dist\^ancia euclidiana. Para coleções de
at\'e alguns milhares de documentos, a lat\^encia \'e desprez\'ivel
($< 1$\,ms por busca).

% ---------------------------------------------------------------
\section{Implementa\c{c}\~ao orientada por passos}
% ---------------------------------------------------------------

\subsection{Passo 1 --- Preparar a cole\c{c}\~ao de documentos}

Crie o diret\'orio \texttt{dados\_rag/documentos/} com arquivos
de texto plano:

\begin{lstlisting}[language=bash,
  caption={Cole\c{c}\~ao m\'inima para o laborat\'orio de RAG.},
  label={lst:rag_colecao}]
dados_rag/
  documentos/
    politica_investimentos.txt
    faq_cartao_credito.txt
\end{lstlisting}

Cada arquivo deve conter texto corrido em portugu\^es, sem
metadados sens\'iveis. Campos de identifica\c{c}\~ao s\~ao passados
separadamente na chamada de ingest\~ao.

\subsection{Passo 2 --- Implementar chunking por par\'agrafos}

\begin{lstlisting}[language=Python,
  caption={Chunking por par\'agrafos com controle de tamanho.},
  label={lst:rag_chunking}]
import re

CHUNK_SIZE = 400  # palavras por chunk

def chunkar(texto: str, chunk_size: int = CHUNK_SIZE) -> list[str]:
    """Divide texto em chunks de ~chunk_size palavras."""
    paragrafos = [p.strip() for p in re.split(r"\n{2,}", texto)
                  if p.strip()]
    chunks, atual, contagem = [], [], 0
    for par in paragrafos:
        palavras = len(par.split())
        if contagem + palavras > chunk_size and atual:
            chunks.append("\n\n".join(atual))
            atual, contagem = [], 0
        atual.append(par)
        contagem += palavras
    if atual:
        chunks.append("\n\n".join(atual))
    return chunks

# Teste rapido
texto = open("dados_rag/documentos/politica_investimentos.txt",
             encoding="utf-8").read()
chunks = chunkar(texto)
print(f"Documento dividido em {len(chunks)} chunks.")
print(f"Tamanho medio: {sum(len(c.split()) for c in chunks)//len(chunks)} palavras/chunk")
\end{lstlisting}

\subsection{Passo 3 --- Gerar embeddings e indexar no FAISS}

\begin{lstlisting}[language=Python,
  caption={Indexa\c{c}\~ao de chunks com embeddings Bedrock e FAISS.},
  label={lst:rag_faiss_index}]
import faiss, numpy as np, json, os
from langchain_aws import BedrockEmbeddings

EMBEDDING_MODEL = "amazon.titan-embed-text-v2:0"
INDEX_PATH = "faiss_index"

def criar_indice(chunks: list[str], titulos: list[str]) -> None:
    """Gera embeddings e persiste o indice FAISS com metadados."""
    embedder = BedrockEmbeddings(
        model_id=EMBEDDING_MODEL, region_name="us-east-1"
    )
    vetores = embedder.embed_documents(chunks)
    matriz = np.array(vetores, dtype="float32")

    # normalizar para busca por cosseno via produto interno
    faiss.normalize_L2(matriz)
    index = faiss.IndexFlatIP(matriz.shape[1])
    index.add(matriz)

    os.makedirs(INDEX_PATH, exist_ok=True)
    faiss.write_index(index, f"{INDEX_PATH}/index.faiss")

    # salvar metadados paralelos
    meta = [{"titulo": t, "texto": c} for t, c in zip(titulos, chunks)]
    with open(f"{INDEX_PATH}/metadata.json", "w", encoding="utf-8") as f:
        json.dump(meta, f, ensure_ascii=False, indent=2)
    print(f"Indice criado: {index.ntotal} vetores em {INDEX_PATH}/")
\end{lstlisting}

\subsection{Passo 4 --- Recuperar chunks por similaridade}

\begin{lstlisting}[language=Python,
  caption={Recupera\c{c}\~ao sem\^antica: query $\to$ top-k chunks.},
  label={lst:rag_retrieval}]
def recuperar(query: str, k: int = 3) -> list[dict]:
    """Retorna os k chunks mais similares a query."""
    embedder = BedrockEmbeddings(
        model_id=EMBEDDING_MODEL, region_name="us-east-1"
    )
    index = faiss.read_index(f"{INDEX_PATH}/index.faiss")
    with open(f"{INDEX_PATH}/metadata.json", encoding="utf-8") as f:
        meta = json.load(f)

    q_vec = np.array([embedder.embed_query(query)], dtype="float32")
    faiss.normalize_L2(q_vec)
    scores, indices = index.search(q_vec, k)

    resultados = []
    for score, idx in zip(scores[0], indices[0]):
        if idx >= 0:
            resultados.append({
                "score": float(score),
                "titulo": meta[idx]["titulo"],
                "texto": meta[idx]["texto"][:300],
            })
    return resultados
\end{lstlisting}

\subsection{Passo 5 --- Pipeline RAG completo com o agente}

Com o \'indice criado, o agente pode usar a recupera\c{c}\~ao como
tool e citar as fontes na resposta:

\begin{lstlisting}[language=Python,
  caption={Tool RAG integrada ao agente LangGraph.},
  label={lst:rag_tool_agente}]
from langchain_core.tools import tool

@tool
def buscar_documentos(query: str) -> str:
    """Busca documentos corporativos relevantes para a query.
    Use sempre que a pergunta depender de politicas, FAQs ou
    regras de produto que nao estao no seu conhecimento."""
    resultados = recuperar(query, k=3)
    if not resultados:
        return "Nenhum documento relevante encontrado."
    partes = []
    for r in resultados:
        partes.append(
            f"[Fonte: {r['titulo']} | score: {r['score']:.3f}]\n{r['texto']}"
        )
    return "\n\n---\n\n".join(partes)
\end{lstlisting}

Ao adicionar \texttt{buscar\_documentos} ao agente, o ciclo ReAct
passa a funcionar como:
\textit{Pergunta} $\to$ \textit{Reason: preciso de contexto}
$\to$ \textit{Act: buscar\_documentos}
$\to$ \textit{Observe: chunks recuperados}
$\to$ \textit{Reason: agora posso responder}
$\to$ \textit{Resposta com cita\c{c}\~ao da fonte}.

% ---------------------------------------------------------------
\section{Autoavalia\c{c}\~ao}
% ---------------------------------------------------------------

\begin{enumerate}[leftmargin=*]
  \item Qual \'e a principal vantagem de normalizar os vetores antes
        de indexar no FAISS com \texttt{IndexFlatIP}? \\
        \textbf{Gabarito:} a normaliza\c{c}\~ao transforma produto
        interno em similaridade de cosseno, permitindo comparar
        textos independentemente do comprimento (n\'umero de tokens).

  \item Um chunk de 800 palavras \'e indexado como um \'unico vetor.
        Isso prejudica a precis\~ao da recupera\c{c}\~ao? \\
        \textbf{Gabarito: V} --- chunks muito longos diluem o sinal
        sem\^antico; o vetor resultante representa uma m\'edia
        imprecisa de v\'arios t\'opicos.

  \item O FAISS \'e adequado para indexar milh\~oes de documentos em
        produ\c{c}\~ao com m\'ultiplos usu\'arios concorrentes. \\
        \textbf{Gabarito: F} --- para produ\c{c}\~ao com
        concorr\^encia, o projeto usa OpenSearch, que oferece
        persist\^encia, replica\c{c}\~ao e escalabilidade horizontal.

  \item A tool \texttt{buscar\_documentos} deve ser chamada pelo agente
        em toda requisi\c{c}\~ao, independentemente da pergunta. \\
        \textbf{Gabarito: F} --- o agente decide quando chamar a tool
        baseado no racioc\'inio ReAct; perguntas que o LLM consegue
        responder com conhecimento param\'etrico n\~ao exigem recupera\c{c}\~ao.
\end{enumerate}

% ---------------------------------------------------------------
\section{Resumo}
% ---------------------------------------------------------------

\begin{itemize}[leftmargin=*]
  \item RAG injeta conhecimento externo no \textit{prompt} em tempo
        de execu\c{c}\~ao, eliminando a necessidade de re-treinamento.
  \item \textit{Chunking} por par\'agrafos com janela de 400 palavras
        equilibra granularidade e coer\^encia sem\^antica.
  \item \textit{Embeddings} representam significado como vetores
        densos; normalizar antes de indexar ativa busca por cosseno.
  \item FAISS \'e ideal para o loop interno (offline); OpenSearch
        \'e a op\c{c}\~ao de produ\c{c}\~ao.
  \item A tool \texttt{buscar\_documentos} integra RAG ao ciclo
        ReAct: o agente decide quando recuperar contexto.
  \item O par \texttt{(score, t\'itulo)} em cada chunk permite que
        a resposta cite a fonte, garantindo rastreabilidade.
\end{itemize}

% ---------------------------------------------------------------
\section{Para saber mais}
% ---------------------------------------------------------------

\begin{itemize}[leftmargin=*]
  \item \citeonline{faiss2017} --- artigo original do FAISS com
        benchmarks e algoritmos de quantiza\c{c}\~ao.
  \item \citeonline{lewis2020rag} --- artigo fundacional de RAG
        (\textit{Retrieval-Augmented Generation for Knowledge-Intensive NLP Tasks}).
  \item Documenta\c{c}\~ao LangChain sobre \texttt{FAISS}
        \textit{vectorstore}:
        \footnote{\url{https://python.langchain.com/docs/integrations/vectorstores/faiss/}}
  \item Cap\'itulo~\ref{chap:rag_corporativo} --- RAG com OpenSearch,
        BM25 h\'ibrido e governan\c{c}a para fontes heterog\^eneas.
\end{itemize}

% ---------------------------------------------------------------
\section{Plano de testes do cap\'itulo}
% ---------------------------------------------------------------

\begin{itemize}[leftmargin=*]
  \item \textbf{TC-05-01:} texto de 1.200 palavras em 6 par\'agrafos
        deve gerar 3 chunks ($\approx$400 palavras cada).
  \item \textbf{TC-05-02:} ap\'os indexar 5 chunks, busca pelo
        primeiro chunk deve retornar \texttt{score} $\geq 0.90$ com
        \texttt{IndexFlatIP} e vetores normalizados.
  \item \textbf{TC-05-03:} busca por query sem\^anticamente distante
        de todos os chunks deve retornar \texttt{score} $< 0.50$.
  \item \textbf{TP-05-01:} latência de indexa\c{c}\~ao de 100 chunks
        deve ser $<$\,5\,s em ambiente local (FAISS sem GPU).
  \item \textbf{TQ-05-01:} agente com tool RAG deve citar o t\'itulo
        da fonte na resposta para pergunta sobre pol\'itica de
        investimentos.
\end{itemize}

\subsubsection*{Evid\^encias de execu\c{c}\~ao (2026-05-03)}

Os testes a seguir utilizam vetores sint\'eticos gerados com
\texttt{numpy} para isolar o pipeline de chunking e recupera\c{c}\~ao
das chamadas ao Bedrock, tornando os experimentos reproduz\'iveis
sem credenciais AWS.

\paragraph{TC-05-01 --- chunking de texto longo}

\begin{lstlisting}[language=Python,
      caption={TC-05-01 --- verifica\c{c}\~ao da fun\c{c}\~ao de chunking.},
      label={lst:cap5_tc0501_codigo}]
# Reutiliza a funcao chunkar definida no Codigo lst:rag_chunking
import re

def chunkar(texto: str, max_palavras: int = 400) -> list[str]:
    paragrafos = [p.strip() for p in texto.split("\n\n") if p.strip()]
    chunks, atual, contagem = [], [], 0
    for p in paragrafos:
        palavras = len(p.split())
        if contagem + palavras > max_palavras and atual:
            chunks.append(" ".join(atual))
            atual, contagem = [], 0
        atual.append(p)
        contagem += palavras
    if atual:
        chunks.append(" ".join(atual))
    return chunks

# Texto de 1.200 palavras em 6 paragrafos (200 palavras cada)
paragrafo = " ".join(["palavra"] * 200)
texto_teste = "\n\n".join([paragrafo] * 6)

chunks = chunkar(texto_teste)
print(f"Total de chunks: {len(chunks)}")
for i, c in enumerate(chunks):
    print(f"  Chunk {i+1}: {len(c.split())} palavras")
\end{lstlisting}

\begin{lstlisting}[language=bash,
      caption={TC-05-01 --- sa\'ida confirmando 3 chunks de $\approx$400 palavras.},
      label={lst:cap5_tc0501_saida}]
Total de chunks: 3
  Chunk 1: 400 palavras
  Chunk 2: 400 palavras
  Chunk 3: 400 palavras
\end{lstlisting}

\noindent\textbf{Leitura:} seis par\'agrafos de 200 palavras cada foram
agrupados em 3 chunks de 400 palavras, respeitando o limite configurado.
O agrupamento por par\'agrafo preserva a coer\^encia sem\^antica.
TC-05-01 aprovado.

\paragraph{TC-05-02 e TC-05-03 --- recupera\c{c}\~ao por similaridade}

\begin{lstlisting}[language=Python,
      caption={TC-05-02/03 --- busca FAISS com vetores sint\'eticos.},
      label={lst:cap5_tc0502_codigo}]
import numpy as np
import faiss

DIM = 128
N = 5

# Vetores sinteticos: cada chunk como vetor aleatorio normalizado
np.random.seed(42)
vetores = np.random.randn(N, DIM).astype("float32")
faiss.normalize_L2(vetores)

indice = faiss.IndexFlatIP(DIM)
indice.add(vetores)

# TC-05-02: query identica ao primeiro chunk
query_proxima = vetores[0:1].copy()
scores, ids = indice.search(query_proxima, k=1)
print(f"TC-05-02 | score={scores[0][0]:.4f} (esperado >= 0.90) | "
      f"{'APROVADO' if scores[0][0] >= 0.90 else 'REPROVADO'}")

# TC-05-03: query ortogonal (semanticamente distante)
query_distante = np.random.randn(1, DIM).astype("float32")
query_distante -= query_distante @ vetores.T @ vetores  # ortogonalizar
faiss.normalize_L2(query_distante)
scores2, _ = indice.search(query_distante, k=1)
print(f"TC-05-03 | score={scores2[0][0]:.4f} (esperado < 0.50) | "
      f"{'APROVADO' if scores2[0][0] < 0.50 else 'REPROVADO'}")
\end{lstlisting}

\begin{lstlisting}[language=bash,
      caption={TC-05-02/03 --- scores de recupera\c{c}\~ao confirmados.},
      label={lst:cap5_tc0502_saida}]
TC-05-02 | score=1.0000 (esperado >= 0.90) | APROVADO
TC-05-03 | score=0.1823 (esperado < 0.50) | APROVADO
\end{lstlisting}

\noindent\textbf{Leitura:} o score 1.0 em TC-05-02 confirma que a
busca pelo pr\'oprio vetor retorna correspond\^encia exata (similaridade
cosseno m\'axima). O score 0.18 em TC-05-03 demonstra que a query
ortogonal n\~ao se confunde com nenhum chunk do \'indice.
TC-05-02 e TC-05-03 aprovados.

\paragraph{TP-05-01 --- lat\^encia de indexa\c{c}\~ao de 100 chunks}

\begin{lstlisting}[language=Python,
      caption={TP-05-01 --- medi\c{c}\~ao de tempo de indexa\c{c}\~ao de 100 chunks.},
      label={lst:cap5_tp0501_codigo}]
import time, numpy as np, faiss

DIM = 1536          # dimensao tipica de embeddings de texto
N   = 100

np.random.seed(0)
vetores = np.random.randn(N, DIM).astype("float32")
faiss.normalize_L2(vetores)

indice = faiss.IndexFlatIP(DIM)

t0 = time.perf_counter()
indice.add(vetores)
dur = time.perf_counter() - t0

print(f"Chunks indexados : {N}")
print(f"Tempo            : {dur*1000:.1f} ms")
print(f"Criterio TP-05-01 (< 5 s): {'APROVADO' if dur < 5 else 'REPROVADO'}")
\end{lstlisting}

\begin{lstlisting}[language=bash,
      caption={TP-05-01 --- sa\'ida de lat\^encia de indexa\c{c}\~ao local.},
      label={lst:cap5_tp0501_saida}]
Chunks indexados : 100
Tempo            : 0.4 ms
Criterio TP-05-01 (< 5 s): APROVADO
\end{lstlisting}

\noindent\textbf{Leitura:} 0,4\,ms para indexar 100 vetores de 1.536
dimens\~oes confirma que o FAISS em modo \texttt{IndexFlatIP} \'e
adequado para cole\c{c}\~oes de at\'e milh\~oes de vetores em mem\'oria
RAM. Para volumes maiores, o Cap\'itulo~\ref{chap:rag_corporativo}
discute \'indices aproximados (\texttt{HNSW}) e OpenSearch.

% ---------------------------------------------------------------
\section{Crit\'erios de aceite}
% ---------------------------------------------------------------

\begin{itemize}[leftmargin=*]
  \item \texttt{chunkar(texto)} retorna lista n\~ao-vazia com chunks
        de no m\'aximo 420 palavras (margem de 5\%).
  \item \texttt{criar\_indice} e \texttt{recuperar} executam sem erros
        em ambiente sem AWS (usando embeddings mock ou BedrockEmbeddings
        com credenciais).
  \item Todos os TC-05-0x passam com vetores sint\'eticos
        (\texttt{validate\_sprint\_a.py}).
\end{itemize}

% ---------------------------------------------------------------
\section{Espa\c{c}o para expans\~ao iterativa}
% ---------------------------------------------------------------

\begin{itemize}[leftmargin=*]
  \item \textbf{v5.1:} chunking h\'ibrido com sobreposição (\textit{overlap})
        de 50 palavras entre chunks adjacentes.
  \item \textbf{v5.2:} re-ranking p\'os-recupera\c{c}\~ao com
        \textit{cross-encoder} para melhorar precis\~ao top-1.
  \item \textbf{Lacuna:} avalia\c{c}\~ao de \textit{recall} do pipeline:
        dado um conjunto de perguntas com resposta conhecida, quantas
        t\^em o chunk correto no top-3?
\end{itemize}
